% BHU PYQ -- Transcribed & Corrected
% Paper: CHEMJ43: Physical Chemistry
% Exam: B.Sc. (Hons) Semester IV Examination 2025-26 (NEP)
% Ref Code: Conf/Even/N/2025-26/07
% Category: Major
% Department: Chemistry

\documentclass[11pt,a4paper]{article}
\usepackage[a4paper,top=2.5cm,bottom=2.5cm,left=2.8cm,right=2.8cm,headheight=14pt]{geometry}
\usepackage{amsmath,amssymb}
\usepackage[T1]{fontenc}
\usepackage[utf8]{inputenc}
\usepackage{lmodern,microtype,fancyhdr,enumitem,hyperref,lastpage,parskip}

\hypersetup{
    colorlinks=true,
    urlcolor=black,
    linkcolor=black,
    pdftitle={CHEMJ43: Physical Chemistry -- BHU Sem IV 2025-26 (NEP)}
}

\pagestyle{fancy}
\fancyhf{}
\renewcommand{\headrulewidth}{0.4pt}
\renewcommand{\footrulewidth}{0.4pt}
\fancyhead[L]{\small \textit{Banaras Hindu University}}
\fancyhead[R]{\small \textit{CHEMJ43: Physical Chemistry (2025-26)}}
\fancyfoot[C]{\small Page \thepage\ of \pageref{LastPage}}

\newcommand{\pts}[1]{\hfill \textit{[#1]}}

\begin{document}

\begin{center}
  {\large\bfseries BANARAS HINDU UNIVERSITY}\\[2pt]
  {\normalsize B.Sc. (Hons) Semester IV Examination 2025-26 (NEP)}\\[6pt]
  \rule{0.8\linewidth}{0.5pt}\\[6pt]
  {\large\bfseries CHEMISTRY}\\[4pt]
  {\large\bfseries Paper Code: CHEMJ43 : Physical Chemistry}\\[2pt]
  {\small Ref. No.: Conf/Even/N/2025-26/07}\\[4pt]
  {\normalsize Time: 2$\frac{1}{2}$ Hours\hfill Maximum Marks: 70}
\end{center}

\rule{\linewidth}{0.4pt}

\smallskip

\noindent \textit{Instructions: Answer five questions including Question 1 which is compulsory. The figures in the right-hand margin indicate marks.}

\bigskip

\noindent \textbf{Question 1.}
\begin{enumerate}[label=(\alph*)]
  \item Define partial molar volume. For a particular composition, explain how the negative, positive and zero partial molar volumes could be identified from the graph representing the variation of partial molar volume. \pts{2}
  \item What are Raoult's law and Henry's law? How can the ideal solution be defined based on these laws? \pts{2}
  \item Explain any two deactivation processes of a photochemically excited molecule. \pts{2}
  \item Write the equation explaining the effect of the inhibitor on the kinetics of the unimolecular surface reaction and explain, if the inhibitor is strongly adsorbed, then how the rate of reaction varies. \pts{2}
  \item Write briefly about the types of radioactive decay with suitable examples. \pts{2}
  \item During the use of radioisotopes, what are the points to be taken into account? \pts{2}
  \item What do you understand by the purity and strength of radioisotopes? \pts{2}
\end{enumerate}

\bigskip

\noindent \textbf{Question 2.}
\begin{enumerate}[label=(\alph*)]
  \item Define fugacity and fugacity coefficient. When real gas behaves like an ideal gas, what would be the value of the fugacity coefficient? \pts{3}
  \item Write the expressions of Gibbs' energy of mixing and entropy of mixing and from these expressions, explain when and how the mixing of ideal gases is spontaneous. \pts{4}
  \item What is meant by chemical potential? How does the chemical potential vary with temperature and pressure? Derive the Gibbs-Duhem equation. \pts{7}
\end{enumerate}

\bigskip

\noindent \textbf{Question 3.}
\begin{enumerate}[label=(\alph*)]
  \item Write the BET equation and explain how the graph needs to be plotted to obtain a straight line. \pts{3}
  \item What do you mean by (i) Quantum yield and (ii) Law of photochemical equivalence? \pts{4}
  \item Starting from the principle that at equilibrium the chemical potential of the pure solvent vapour equals the chemical potential of the solvent in the liquid mixture, derive the mathematical expression for the elevation of boiling point: $\Delta T = K x_B$. State the assumptions made regarding the solute during this derivation. \pts{7}
\end{enumerate}

\bigskip

\noindent \textbf{Question 4.}
\begin{enumerate}[label=(\alph*)]
  \item Define a photostationary state and explain how it differs fundamentally from a true chemical thermal equilibrium. \pts{3}
  \item Write the rate expression of a unimolecular surface reaction, sketching the graph for the surface covered vs. pressure and explain how the reaction rate varies from first-order at very low pressures to zero-order at extremely high pressures. \pts{4}
  \item Considering a suitable example of fluorescence quenching, derive the Stern-Volmer equation and find out the slope and intercept from the graph between $I_0/I_f$ and $[Q]$. \pts{7}
\end{enumerate}

\bigskip

\noindent \textbf{Question 5.}
\begin{enumerate}[label=(\alph*)]
  \item Discuss the influence of pH and temperature, depicting the graph, on the kinetics of enzyme-catalysed reactions. \pts{3}
  \item Discuss the photochemical decomposition of HI and find out its quantum yield. \pts{4}
  \item Using the steady-state principle for an intermediate enzyme-substrate complex (ES), derive the Michaelis-Menten equation. Draw the Lineweaver-Burk plot and explain what value could be obtained from its slope and intercept. \pts{7}
\end{enumerate}

\bigskip

\noindent \textbf{Question 6.}
\begin{enumerate}[label=(\alph*)]
  \item Explain, with examples, transfer and reverse transfer reactions. \pts{3}
  \item Calculate the excitation energy of the compound nucleus $^{26}\text{Al}^*$ formed by the bombardment of $^{24}\text{Mg}$ with an $8.0\text{ MeV}$ deuteron ($d$). [Given: $m(^{24}\text{Mg}) = 23.9850\text{ u}$, $m(d) = 2.0141\text{ u}$, $m(^{26}\text{Al}) = 25.9869\text{ u}$ and K.E. of projectile and target on Center of Mass System (CMS) = $7.38\text{ MeV}$]. \pts{3}
  \item What is artificial radioactivity? Describe the experimental setup and evidence used by Frédéric and Irène Joliot-Curie to demonstrate the discovery of artificial radioactivity. \pts{4}
  \item Prove that the kinetic energy available for excitation of the target in the Lab System (LS) is equal to the total kinetic energy of the projectile and target in the CMS. \pts{4}
\end{enumerate}

\bigskip

\noindent \textbf{Question 7.}
\begin{enumerate}[label=(\alph*)]
  \item How can one determine the structure of $\text{PCl}_5$ using radioisotopes as tracers? \pts{3}
  \item Explain the principle of Isotope Dilution Analysis for determining the total blood volume in a human body. \pts{3}
  \item Define the nuclear reaction cross-section. Derive the relationship between the cross-section and the reaction rate for a beam of particles incident on a target. \pts{4}
  \item What is a Fricke dosimeter? Explain the detailed mechanism of the radiolysis in the Fricke dosimeter aqueous solution. \pts{4}
\end{enumerate}

\bigskip

\noindent \textbf{Question 8.} Answer any two of the following: \pts{2 $\times$ 7 = 14}
\begin{enumerate}[label=(\alph*)]
  \item Write the different steps for the photochemical reaction between $\text{H}_2$ and $\text{Br}_2$, and obtain the expression for the rate of formation of $\text{HBr}$.
  \item What is Osmometry? Find out the osmotic pressure at $25\ ^\circ\text{C}$ of a solution containing $1\text{ g}$ glucose ($\text{M.Wt. } 180\text{ g/mole}$) and $1\text{ g}$ sucrose ($\text{M.Wt. } 342\text{ g/mole}$) per litre.
  \item Write about Compton Scattering.
\end{enumerate}

\bigskip
\begin{center}
  \textbf{***}
\end{center}

\end{document}
